% =====================================================================
% Chủ đề 4 — Số học và Tư duy số — phần c: bài 31–40 (mở rộng từ bài 1–15)
% =====================================================================
\clearpage
\dividerbar{Mở rộng -- Nâng cao}

% ---------------------------------------------------------------- Bài 31
\needspace{9\baselineskip}
\begin{probbox}{31}{\relax}
Cho các số: $-1{,}5;\quad 0{,}25;\quad -0{,}75;\quad 1{,}2;\quad -2$.
\begin{enumerate}[label=\alph*),leftmargin=1.6em,itemsep=0pt,topsep=2pt]
  \item Sắp xếp các số theo thứ tự tăng dần.
  \item Những số nào nằm giữa $-1$ và $1$ trên trục số?
  \item Tính tổng của số lớn nhất và số nhỏ nhất.
\end{enumerate}
\end{probbox}
\needspace{7\baselineskip}
\probsol{
\textbf{Phân tích đề.} Bài này kết hợp hai kỹ năng của Bài 1 và Bài 3: so sánh số thập phân và so sánh số âm.
\begin{recallbox}Số âm luôn nhỏ hơn số dương. Trong hai số âm, số nào xa $0$ hơn thì nhỏ hơn. Trong hai số dương,
so sánh phần nguyên rồi từng hàng thập phân.\end{recallbox}
\step{1}{Chia thành hai nhóm: số âm và số dương}
Số âm: $-1{,}5;\ -0{,}75;\ -2$. \quad Số dương: $0{,}25;\ 1{,}2$.
\step{2}{Sắp xếp nhóm số âm}
Khoảng cách đến $0$: $-2$ cách $2$; \ $-1{,}5$ cách $1{,}5$; \ $-0{,}75$ cách $0{,}75$. Số xa $0$ nhất là nhỏ nhất:
\[ -2 < -1{,}5 < -0{,}75. \]
\step{3}{Sắp xếp nhóm số dương rồi ghép lại}
$0{,}25 < 1{,}2$ (phần nguyên $0 < 1$). Ghép: số âm đứng trước, số dương đứng sau:
\[ -2 < -1{,}5 < -0{,}75 < 0{,}25 < 1{,}2. \]
\begin{figC}\begin{scope}[x=1.6cm]
  \draw[tmGreen!35,line width=3pt] (-1,0)--(1,0);
  \draw[->,thick,Navy] (-2.5,0)--(1.7,0);
  \foreach \x in {-2,-1,0,1}{\draw[Navy] (\x,0.08)--(\x,-0.08);}
  \foreach \x in {-1,0,1}{\node[below,font=\scriptsize] at (\x,-0.1) {$\x$};}
  \foreach \x/\l in {-2/{-2},-1.5/{-1{,}5},-0.75/{-0{,}75},0.25/{0{,}25},1.2/{1{,}2}}{\fill[tmOrange] (\x,0) circle (2.6pt); \node[above,font=\scriptsize] at (\x,0.1) {$\l$};}
\end{scope}\end{figC}
\step{4}{Các số nằm giữa $-1$ và $1$}
Trên trục số, đoạn tô xanh từ $-1$ đến $1$ chứa $-0{,}75$ và $0{,}25$. (Số $-1{,}5$ ở bên trái $-1$; số $1{,}2$ ở bên phải $1$.)
\step{5}{Tổng của số lớn nhất và số nhỏ nhất}
Số lớn nhất là $1{,}2$, số nhỏ nhất là $-2$:
\[ 1{,}2 + (-2) = -(2 - 1{,}2) = -0{,}8. \]
\begin{recallbox}Cộng hai số trái dấu: lấy số lớn trừ số bé (không kể dấu), rồi đặt dấu của số có khoảng cách đến $0$ lớn hơn.\end{recallbox}
\begin{warnbox}Nhiều bạn xếp $-0{,}75$ trước $-1{,}5$ vì ``$0{,}75 < 1{,}5$''. Với số âm thì ngược lại: $-1{,}5 < -0{,}75$.\end{warnbox}
\verify{$-0{,}8 - 1{,}2 = -2$. Đúng.}
\ans{a) $-2 < -1{,}5 < -0{,}75 < 0{,}25 < 1{,}2$;\quad b) $-0{,}75$ và $0{,}25$;\quad c) $-0{,}8$}
}

% ---------------------------------------------------------------- Bài 32
\needspace{9\baselineskip}
\begin{probbox}{32}{\relax}
Cho biểu thức $2 + 3 \times 4 - 1$.
\begin{enumerate}[label=\alph*),leftmargin=1.6em,itemsep=0pt,topsep=2pt]
  \item Tính giá trị của biểu thức.
  \item Đặt thêm \emph{một} cặp dấu ngoặc $(\ )$ vào biểu thức (không đổi thứ tự các số và các dấu) để được giá trị
  lớn nhất có thể. Giá trị lớn nhất đó là bao nhiêu?
  \item Cũng hỏi như câu b) nhưng với giá trị nhỏ nhất.
\end{enumerate}
\end{probbox}
\needspace{7\baselineskip}
\probsol{
\textbf{Phân tích đề.} Bài này mở rộng Bài 4 và Bài 7: vị trí dấu ngoặc làm thay đổi thứ tự tính, nên làm thay đổi kết quả.
\begin{recallbox}Trong ngoặc trước; rồi nhân, chia; cuối cùng cộng, trừ (từ trái sang phải).\end{recallbox}
\step{1}{Giá trị khi chưa có ngoặc}
$2 + 3\times 4 - 1 = 2 + 12 - 1 = 13$.
\step{2}{Liệt kê có hệ thống mọi cách đặt một cặp ngoặc}
Một cặp ngoặc phải bao ít nhất hai số liền nhau. Gọi bốn số theo thứ tự là số thứ $1, 2, 3, 4$ ($2;\ 3;\ 4;\ 1$).
Ngoặc có thể bắt đầu ở số thứ $i$ và kết thúc ở số thứ $j$ với $i < j$: có $6$ cách.
\begin{center}
\begin{tabular}{|c|l|c|}
\hline
\rowcolor{tmBlue!10} Ngoặc bao & Biểu thức & Giá trị\\ \hline
số 1 -- 2 & $(2+3)\times 4 - 1 = 5\times 4 - 1$ & $\mathbf{19}$\\ \hline
số 1 -- 3 & $(2+3\times 4) - 1 = 14 - 1$ & $13$\\ \hline
số 1 -- 4 & $(2+3\times 4-1)$ & $13$\\ \hline
số 2 -- 3 & $2 + (3\times 4) - 1$ & $13$\\ \hline
số 2 -- 4 & $2 + (3\times 4 - 1) = 2 + 11$ & $13$\\ \hline
số 3 -- 4 & $2 + 3\times(4 - 1) = 2 + 9$ & $\mathbf{11}$\\ \hline
\end{tabular}
\end{center}
\step{3}{Đọc kết quả}
Lớn nhất là $19$, khi viết $(2+3)\times 4 - 1$: ta ``phóng to'' tổng $2 + 3$ lên $4$ lần.\\
Nhỏ nhất là $11$, khi viết $2 + 3\times(4-1)$: ta nhân $3$ với một số nhỏ hơn ($3$ thay vì $4$).
\begin{tipbox}Muốn kết quả lớn, hãy để phép nhân tác động lên một tổng lớn; muốn kết quả nhỏ, hãy để phép nhân tác động
lên một hiệu. Nhưng để chắc chắn, cách tốt nhất vẫn là liệt kê đủ mọi trường hợp như trong bảng.\end{tipbox}
\verify{$(2+3)\times 4 - 1 = 20 - 1 = 19$ và $2 + 3\times 3 = 11$.}
\ans{a) $13$;\quad b) $(2+3)\times 4 - 1 = 19$;\quad c) $2 + 3\times(4-1) = 11$}
}

% ---------------------------------------------------------------- Bài 33
\needspace{9\baselineskip}
\begin{probbox}{33}{\relax}
Cuốn truyện của Anna dày $360$ trang. Ngày thứ nhất Anna đọc $\dfrac{1}{4}$ cuốn truyện. Ngày thứ hai Anna đọc
$\dfrac{1}{3}$ \emph{số trang còn lại}. Hỏi sau hai ngày, Anna còn bao nhiêu trang chưa đọc, và số trang đó bằng mấy phần
cuốn truyện?
\end{probbox}
\needspace{7\baselineskip}
\probsol{
\textbf{Phân tích đề.} Cần chú ý: phân số $\dfrac13$ của ngày thứ hai được tính trên \emph{số trang còn lại}, không phải
trên cả cuốn truyện. Đây là phiên bản ``hai bước'' của Bài 9.
\begin{recallbox}$\dfrac{m}{n}$ của $a$ bằng $a : n \times m$.\end{recallbox}
\step{1}{Ngày thứ nhất}
Anna đọc $360 : 4 = 90$ (trang). Còn lại $360 - 90 = 270$ (trang).
\step{2}{Ngày thứ hai (tính trên $270$ trang)}
Anna đọc $270 : 3 = 90$ (trang). Còn lại $270 - 90 = 180$ (trang).
\step{3}{Phần cuốn truyện còn lại}
Số trang chưa đọc chiếm $\dfrac{180}{360} = \dfrac{180:180}{360:180} = \dfrac12$ cuốn truyện.
\begin{warnbox}Sai lầm phổ biến: cộng $\dfrac14 + \dfrac13 = \dfrac{7}{12}$ rồi tính phần còn lại là $\dfrac{5}{12}$ cuốn
($150$ trang). Sai, vì $\dfrac13$ là của \emph{phần còn lại} ($270$ trang), không phải của cả cuốn ($360$ trang).\end{warnbox}
\step{4}{Cách khác: nhân các phân số ``còn lại''}
Sau ngày 1 còn $\dfrac34$ cuốn; sau ngày 2 còn $\dfrac23$ của phần đó, tức $\dfrac23\times\dfrac34 = \dfrac{6}{12} = \dfrac12$ cuốn,
bằng $360 : 2 = 180$ trang.
\verify{$90 + 90 + 180 = 360$. Đúng.}
\ans{$180$ trang, bằng $\dfrac12$ cuốn truyện}
}

% ---------------------------------------------------------------- Bài 34
\needspace{9\baselineskip}
\begin{probbox}{34}{\relax}
Bài kiểm tra thứ nhất có $40$ câu hỏi, Ben làm đúng $32$ câu.
\begin{enumerate}[label=\alph*),leftmargin=1.6em,itemsep=0pt,topsep=2pt]
  \item Ben làm đúng bao nhiêu phần trăm số câu của bài thứ nhất?
  \item Bài kiểm tra thứ hai có $60$ câu hỏi. Ben phải làm đúng ít nhất bao nhiêu câu ở bài thứ hai để tính chung cả
  hai bài, số câu đúng đạt ít nhất $85\%$ tổng số câu?
\end{enumerate}
\end{probbox}
\needspace{7\baselineskip}
\probsol{
\textbf{Phân tích đề.} Câu a) là tìm tỉ số phần trăm (ngược với Bài 5). Câu b) là tìm phần trăm của một tổng.
\begin{recallbox}Tỉ số phần trăm của $a$ so với $b$ là $\dfrac{a}{b}\times 100\%$. \ $p\%$ của $a$ là $a\times\dfrac{p}{100}$.\end{recallbox}
\step{1}{Tỉ số phần trăm của bài thứ nhất}
$\dfrac{32}{40}\times 100\% = 0{,}8\times 100\% = 80\%$.
\step{2}{Tổng số câu của cả hai bài}
$40 + 60 = 100$ (câu).
\step{3}{Số câu đúng cần có tính chung}
$85\%$ của $100$ câu là $100\times\dfrac{85}{100} = 85$ (câu).
\step{4}{Số câu đúng cần ở bài thứ hai}
Ben đã có $32$ câu đúng, nên cần thêm ít nhất $85 - 32 = 53$ (câu) ở bài thứ hai.
\begin{warnbox}Không được lấy ``trung bình của hai tỉ lệ phần trăm'': nếu tìm $x$ sao cho $(80\% + x) : 2 = 85\%$ thì được
$x = 90\%$ của $60$ câu là $54$ câu -- sai! Hai bài có số câu khác nhau, nên phải tính trên \emph{tổng số câu}.\end{warnbox}
\verify{Với $53$ câu đúng: $(32 + 53) : 100 = 85\%$ (vừa đủ). Với $52$ câu: $84 : 100 = 84\% < 85\%$ (không đủ).}
\ans{a) $80\%$;\quad b) ít nhất $53$ câu}
}

% ---------------------------------------------------------------- Bài 35
\needspace{9\baselineskip}
\begin{probbox}{35}{\relax}
Một cửa hàng bán hạt điều đóng gói: gói $3$ kg giá \$45, gói $5$ kg giá \$70 (không bán lẻ).
\begin{enumerate}[label=\alph*),leftmargin=1.6em,itemsep=0pt,topsep=2pt]
  \item Loại gói nào rẻ hơn tính theo mỗi ki-lô-gam?
  \item Cô Grace cần mua \emph{đúng} $15$ kg. Cô nên mua thế nào để tốn ít tiền nhất? Hết bao nhiêu tiền?
  \item Nếu cần mua \emph{đúng} $13$ kg thì hết bao nhiêu tiền?
\end{enumerate}
\end{probbox}
\needspace{7\baselineskip}
\probsol{
\textbf{Phân tích đề.} Mở rộng Bài 6: so sánh đơn giá, rồi tìm cách ghép các gói cho đúng khối lượng.
\begin{recallbox}Đơn giá $=$ thành tiền $:$ số lượng.\end{recallbox}
\step{1}{Đơn giá của từng loại gói}
Gói $3$ kg: $45 : 3 = 15$ (đô la/kg). \quad Gói $5$ kg: $70 : 5 = 14$ (đô la/kg).\\
Vậy gói $5$ kg rẻ hơn ($14 < 15$).
\step{2}{Các cách ghép đúng $15$ kg}
Gọi $a$ là số gói $3$ kg, $b$ là số gói $5$ kg ($a$, $b$ là số tự nhiên). Ta cần $3a + 5b = 15$, tức $3a = 15 - 5b$.
Vì $5b \le 15$ nên $b \in \{0;1;2;3\}$. Thử từng giá trị:
\begin{center}
\begin{tabular}{|c|c|c|c|}
\hline
\rowcolor{tmBlue!10} $b$ & $15 - 5b$ & $a$ & Thành tiền (đô la)\\ \hline
$0$ & $15$ & $5$ & $5\times 45 = 225$\\ \hline
$1$ & $10$ & không chia hết cho $3$ & --\\ \hline
$2$ & $5$ & không chia hết cho $3$ & --\\ \hline
$3$ & $0$ & $0$ & $3\times 70 = \mathbf{210}$\\ \hline
\end{tabular}
\end{center}
Mua $3$ gói $5$ kg là rẻ nhất: \$210.
\step{3}{Đúng $13$ kg}
$3a + 5b = 13$, với $5b \le 13$ nên $b \in \{0;1;2\}$: \ $b = 0 \Rightarrow 3a = 13$ (không chia hết cho $3$); \ $b = 1 \Rightarrow 3a = 8$ (không chia hết cho $3$);
\ $b = 2 \Rightarrow 3a = 3$, $a = 1$. Chỉ có một cách: $1$ gói $3$ kg và $2$ gói $5$ kg.\\
Thành tiền: $45 + 2\times 70 = 45 + 140 = 185$ (đô la).
\begin{warnbox}Không thể trả lời câu c) bằng cách lấy $13\times 14 = 182$ đô la, vì cửa hàng không bán lẻ: phải ghép các gói sao cho
đúng $13$ kg.\end{warnbox}
\verify{$3 + 5 + 5 = 13$ kg và $45 + 70 + 70 = 185$; \ $5 + 5 + 5 = 15$ kg và $70 + 70 + 70 = 210 < 225$.}
\ans{a) gói $5$ kg (\$14/kg);\quad b) $3$ gói $5$ kg, \$210;\quad c) \$185}
}

% ---------------------------------------------------------------- Bài 36
\needspace{9\baselineskip}
\begin{probbox}{36}{\relax}
Năm bạn trong một nhóm có cân nặng (kg) là $38;\ 41;\ 40;\ 37$ và $x$. Cân nặng trung bình của nhóm là $40$ kg.
\begin{enumerate}[label=\alph*),leftmargin=1.6em,itemsep=0pt,topsep=2pt]
  \item Tìm $x$. Bạn thứ năm nặng hơn hay nhẹ hơn mức trung bình bao nhiêu ki-lô-gam?
  \item Tìm trung vị cân nặng của nhóm.
  \item Một bạn thứ sáu gia nhập nhóm, và cân nặng trung bình của cả sáu bạn là $41$ kg. Bạn thứ sáu nặng bao nhiêu?
\end{enumerate}
\end{probbox}
\needspace{7\baselineskip}
\probsol{
\textbf{Phân tích đề.} Mở rộng Bài 14: dùng tổng và độ lệch để tìm số còn thiếu; thêm trung vị và bài toán ``thêm một người''.
\begin{recallbox}Tổng $=$ trung bình cộng $\times$ số các số. Tổng các độ lệch so với trung bình cộng bằng $0$.\end{recallbox}
\step{1}{Tìm $x$ bằng độ lệch}
So với $40$: $38 \to -2$; \ $41 \to +1$; \ $40 \to 0$; \ $37 \to -3$. Tổng: $-2 + 1 + 0 - 3 = -4$.\\
Để tổng các độ lệch bằng $0$, bạn thứ năm phải lệch $+4$: $x = 40 + 4 = 44$ (kg), nặng hơn mức trung bình $4$ kg.\\
Kiểm tra bằng tổng: $5\times 40 = 200$; \ $200 - (38 + 41 + 40 + 37) = 200 - 156 = 44$.
\step{2}{Trung vị}
Sắp xếp tăng dần: $37;\ 38;\ \mathbf{40};\ 41;\ 44$. Số đứng giữa (thứ ba) là $40$. Trung vị là $40$ kg.
\step{3}{Bạn thứ sáu}
Tổng cân nặng sáu bạn: $6\times 41 = 246$ (kg). Năm bạn đầu nặng tổng cộng $200$ kg. Bạn thứ sáu nặng $246 - 200 = 46$ (kg).
\begin{tipbox}Có thể nghĩ bằng độ lệch: trung bình tăng $1$ kg cho cả $6$ người, tức cần thêm $6$ kg so với mức $40$;
bạn mới phải nặng $40 + 6 = 46$ kg.\end{tipbox}
\begin{warnbox}Không phải ``bạn thứ sáu nặng $41$ kg'' -- nếu thế, trung bình chỉ tăng một chút: $241 : 6 \approx 40{,}17$ kg.\end{warnbox}
\verify{$(37 + 38 + 40 + 41 + 44 + 46) : 6 = 246 : 6 = 41$. Đúng.}
\ans{a) $x = 44$ kg, nặng hơn trung bình $4$ kg;\quad b) $40$ kg;\quad c) $46$ kg}
}

% ---------------------------------------------------------------- Bài 37
\needspace{9\baselineskip}
\begin{probbox}{37}{\stars{1}}
Một đường chạy vòng quanh sân trường dài $200$ m. Mia chạy $6$ vòng hết $4$ phút. Tom chạy đều với vận tốc $18$ km/h.
\begin{enumerate}[label=\alph*),leftmargin=1.6em,itemsep=0pt,topsep=2pt]
  \item Ai chạy nhanh hơn?
  \item Nếu giữ nguyên vận tốc, mỗi bạn chạy $3$ km thì hết bao nhiêu phút?
  \item Hai bạn xuất phát cùng lúc, từ cùng một điểm trên đường chạy, nhưng chạy \emph{ngược chiều} nhau. Sau bao nhiêu
  giây thì hai bạn gặp nhau lần đầu tiên?
  \item Một hôm Tom tăng vận tốc lên $21{,}6$ km/h, còn Mia vẫn chạy như cũ. Hai bạn xuất phát cùng lúc, cùng một điểm, chạy
  \emph{cùng chiều}. Sau bao lâu thì Tom vượt Mia đúng một vòng?
\end{enumerate}
\end{probbox}
\needspace{7\baselineskip}
\probsol{
\textbf{Phân tích đề.} Mở rộng Bài 12: tính vận tốc từ số vòng chạy, đổi km/h sang m/s, rồi xét hai người chuyển động
trên cùng một đường vòng.
\begin{recallbox}$v = s : t$, \ $t = s : v$. \ $1$ km/h $= \dfrac{5}{18}$ m/s, tức muốn đổi km/h sang m/s ta nhân với $5$ rồi chia cho $18$.\end{recallbox}
\step{1}{Vận tốc của Mia}
Quãng đường: $6\times 200 = 1200$ (m). Thời gian: $4$ phút $= 240$ giây. \ $v_{\text{Mia}} = 1200 : 240 = 5$ (m/s).
\step{2}{Vận tốc của Tom (đổi ra m/s)}
$v_{\text{Tom}} = 18\times\dfrac{5}{18} = 5$ (m/s). Hai bạn chạy \textbf{nhanh như nhau}!
\begin{warnbox}So sánh trực tiếp ``$5$'' với ``$18$'' là vô nghĩa vì khác đơn vị (m/s và km/h). Phải đổi về cùng đơn vị
trước.\end{warnbox}
\step{3}{Thời gian chạy $3$ km}
$3$ km $= 3000$ m. \ $t = 3000 : 5 = 600$ (giây) $= 600 : 60 = 10$ (phút). Cả hai bạn đều mất $10$ phút.
\step{4}{Chạy ngược chiều}
Mỗi giây Mia chạy $5$ m, Tom chạy $5$ m theo hướng ngược lại, nên quãng đường ``chung'' hai bạn chạy được tăng
$5 + 5 = 10$ m mỗi giây. Hai bạn gặp nhau lần đầu khi \emph{tổng} quãng đường của hai bạn bằng đúng một vòng $200$ m:
\[ t = 200 : 10 = 20 \text{ (giây)}. \]
\begin{figC}
  \draw[Slate!60,line width=4pt] (0,0) circle (1.1);
  \fill[Navy] (0,1.1) circle (2.5pt); \node[above,font=\scriptsize] at (0,1.5) {xuất phát};
  \draw[->,tmBlue,thick] (95:1.4) arc (95:265:1.4);
  \draw[->,tmOrange,thick] (85:1.4) arc (85:-85:1.4);
  \node[left,font=\scriptsize,text=tmBlue!80!black] at (-1.5,0) {Mia: $100$ m};
  \node[right,font=\scriptsize,text=tmOrange!80!black] at (1.5,0) {Tom: $100$ m};
  \fill[AnsRed] (0,-1.1) circle (2.5pt); \node[below,font=\scriptsize,text=AnsRed] at (0,-1.5) {gặp nhau sau $20$ s};
\end{figC}
Sau $20$ giây, mỗi bạn chạy $5\times 20 = 100$ m, đúng nửa vòng: họ gặp nhau ở điểm đối diện với điểm xuất phát.
\step{5}{Chạy cùng chiều, Tom nhanh hơn}
$v_{\text{Tom}} = 21{,}6\times\dfrac{5}{18} = \dfrac{108}{18} = 6$ (m/s). Mỗi giây, Tom chạy hơn Mia $6 - 5 = 1$ m. Tom vượt Mia
đúng một vòng khi Tom chạy hơn Mia $200$ m:
\[ t = 200 : 1 = 200 \text{ (giây)} = 3 \text{ phút } 20 \text{ giây}. \]
\verify{Sau $200$ giây: Tom chạy $6\times 200 = 1200$ m ($6$ vòng), Mia chạy $5\times 200 = 1000$ m ($5$ vòng). Tom hơn đúng $1$ vòng.}
\begin{parentbox}Có thể diễn lại câu c) và d) bằng hai đồng xu trên một chiếc đĩa tròn: ngược chiều thì ``khoảng cách cần
xóa'' giảm theo \emph{tổng} hai vận tốc; cùng chiều thì giảm theo \emph{hiệu} hai vận tốc.\end{parentbox}
\ans{a) nhanh như nhau ($5$ m/s);\quad b) $10$ phút;\quad c) $20$ giây;\quad d) $200$ giây ($3$ phút $20$ giây)}
}

% ---------------------------------------------------------------- Bài 38
\needspace{9\baselineskip}
\begin{probbox}{38}{\stars{1}}
\begin{enumerate}[label=\alph*),leftmargin=1.6em,itemsep=0pt,topsep=0pt]
  \item Tìm các chữ số $a$, $b$ để số $\nd{4a7b}$ chia hết cho cả $2$, $5$ và $9$.
  \item Có bao nhiêu số có bốn chữ số dạng $\nd{5a3b}$ chia hết cho $9$? Hãy liệt kê chúng.
\end{enumerate}
\end{probbox}
\needspace{7\baselineskip}
\probsol{
\textbf{Phân tích đề.} Mở rộng Bài 10: dùng đồng thời nhiều dấu hiệu chia hết và đếm có hệ thống.
\begin{recallbox}Chia hết cho $2$ và $5$ $\iff$ chữ số tận cùng là $0$. \ Chia hết cho $9$ $\iff$ tổng các chữ số chia hết cho $9$.\end{recallbox}
\begin{proofbox}Vì sao ``chia hết cho cả $2$ và $5$'' lại tương đương với ``tận cùng là $0$''? Chia hết cho $2$ cần chữ số
tận cùng chẵn; chia hết cho $5$ cần chữ số tận cùng là $0$ hoặc $5$. Chữ số vừa chẵn vừa thuộc $\{0;5\}$ chỉ có $0$.\end{proofbox}
\step{1}{Câu a) -- tìm $b$}
$\nd{4a7b}$ chia hết cho $2$ và $5$ nên $b = 0$.
\step{2}{Câu a) -- tìm $a$}
Tổng các chữ số: $4 + a + 7 + 0 = 11 + a$. Vì $0 \le a \le 9$ nên $11 \le 11 + a \le 20$. Số chia hết cho $9$ trong khoảng
này chỉ có $18$. Vậy $11 + a = 18$, $a = 7$. Số cần tìm: $4770$.
\verify{$4770 = 9\times 530$, tận cùng là $0$ nên chia hết cho $2$ và $5$.}
\step{3}{Câu b) -- điều kiện}
Tổng các chữ số của $\nd{5a3b}$ là $5 + a + 3 + b = 8 + a + b$. Vì $0 \le a + b \le 18$ nên $8 \le 8 + a + b \le 26$. Các bội
của $9$ trong khoảng này là $9$ và $18$. Do đó $a + b = 1$ hoặc $a + b = 10$.
\step{4}{Câu b) -- đếm có hệ thống}
$a + b = 1$: $(a;b) = (0;1)$ hoặc $(1;0)$ \ -- \ $2$ số.\\
$a + b = 10$: $a$ chạy từ $1$ đến $9$, $b = 10 - a$ (từ $9$ xuống $1$) \ -- \ $9$ số. ($a = 0$ thì $b = 10$, không phải chữ số.)\\
Tổng cộng: $2 + 9 = 11$ số:
\[ 5031,\ 5130,\ 5139,\ 5238,\ 5337,\ 5436,\ 5535,\ 5634,\ 5733,\ 5832,\ 5931. \]
\begin{warnbox}Đừng quên $a$ được phép bằng $0$ (vì chữ số đầu tiên là $5$), và đừng bỏ sót trường hợp tổng chữ số bằng $9$
(tức $a + b = 1$).\end{warnbox}
\verify{Mỗi số trên có tổng chữ số bằng $9$ hoặc $18$, ví dụ $5\,436$: $5+4+3+6 = 18$.}
\ans{a) $a = 7$, $b = 0$ (số $4770$);\quad b) $11$ số}
}

% ---------------------------------------------------------------- Bài 39
\needspace{9\baselineskip}
\begin{probbox}{39}{\stars{1}}
Trong phép cộng dưới đây, $M$ và $N$ là hai chữ số khác nhau và khác $0$. Tìm tất cả các cặp $(M;N)$ thỏa mãn.
\[ \begin{array}{c@{\;\;}c@{\;\;}c@{\;\;}c} & & M & N\\ & & N & M\\ + & & N & N\\ \hline & 1 & 8 & 7 \end{array} \]
\end{probbox}
\needspace{7\baselineskip}
\probsol{
\textbf{Phân tích đề.} Mở rộng Bài 13. Lần này thay vì thử từng hàng, ta dùng cấu tạo số để đưa về một đẳng thức đơn giản,
rồi tìm \emph{tất cả} các nghiệm.
\begin{recallbox}$\nd{MN} = 10M + N$, \ $\nd{NM} = 10N + M$, \ $\nd{NN} = 10N + N = 11N$.\end{recallbox}
\step{1}{Viết tổng bằng cấu tạo số}
Áp dụng cấu tạo số cho ba số hạng, ta có:
\[ \nd{MN} + \nd{NM} + \nd{NN} = (10M + N) + (10N + M) + 11N = 11M + 22N = 11\times(M + 2N). \]
\step{2}{Lập đẳng thức}
$11\times(M + 2N) = 187$, nên $M + 2N = 187 : 11 = 17$.
\step{3}{Tìm mọi cặp chữ số}
$M = 17 - 2N$ phải là chữ số từ $1$ đến $9$, nên $17 - 2N \le 9$, tức $N \ge 4$; và $17 - 2N \ge 1$, tức $N \le 8$.
\begin{center}
\begin{tabular}{|c|c|c|c|c|c|}
\hline
\rowcolor{tmBlue!10} $N$ & $4$ & $5$ & $6$ & $7$ & $8$\\ \hline
$M = 17 - 2N$ & $9$ & $7$ & $5$ & $3$ & $1$\\ \hline
Phép cộng & $94+49+44$ & $75+57+55$ & $56+65+66$ & $37+73+77$ & $18+81+88$\\ \hline
\end{tabular}
\end{center}
Cả năm cặp đều có $M \ne N$, và mỗi tổng đều bằng $187$.
\begin{tipbox}Khi các chữ cái xuất hiện ``đối xứng'' như $\nd{MN}$ và $\nd{NM}$, hãy thử cấu tạo số: thường có một thừa số
chung ($11$ ở đây) giúp bài toán gọn hẳn lại.\end{tipbox}
\begin{warnbox}Đề hỏi \emph{tất cả} các cặp. Tìm được một cặp (chẳng hạn $M = 9$, $N = 4$) rồi dừng là chưa đủ.\end{warnbox}
\verify{$94 + 49 + 44 = 187$; \ $75 + 57 + 55 = 187$; \ $56 + 65 + 66 = 187$; \ $37 + 73 + 77 = 187$; \ $18 + 81 + 88 = 187$.}
\ans{$(M;N) \in \{(9;4),\ (7;5),\ (5;6),\ (3;7),\ (1;8)\}$}
}

% ---------------------------------------------------------------- Bài 40
\needspace{9\baselineskip}
\begin{probbox}{40}{\stars{2}}
Một dãy số được tạo theo quy luật Fibonacci: kể từ số thứ ba, mỗi số bằng tổng của hai số đứng ngay trước nó.
\begin{enumerate}[label=\alph*),leftmargin=1.6em,itemsep=0pt,topsep=2pt]
  \item Biết số thứ nhất là $2$ và số thứ năm là $19$. Tìm số thứ hai và số thứ tám của dãy.
  \item Tính trung bình cộng của năm số đầu tiên của dãy ở câu a).
  \item Trong dãy Fibonacci quen thuộc $1, 1, 2, 3, 5, 8, \ldots$, số thứ $2026$ là số chẵn hay số lẻ? Giải thích.
\end{enumerate}
\end{probbox}
\needspace{7\baselineskip}
\probsol{
\textbf{Phân tích đề.} Mở rộng Bài 15. Câu a) chưa biết số thứ hai, ta gọi nó là $b$ và biểu diễn các số sau theo $b$.
Câu c) không thể viết ra $2026$ số, ta cần tìm quy luật chẵn -- lẻ.
\step{1}{Biểu diễn các số theo $b$}
\begin{center}
\begin{tabular}{|c|c|c|c|c|c|}
\hline
\rowcolor{tmBlue!10} Thứ tự & 1 & 2 & 3 & 4 & 5\\ \hline
Số & $2$ & $b$ & $2 + b$ & $b + (2 + b) = 2 + 2b$ & $(2 + b) + (2 + 2b) = 4 + 3b$\\ \hline
\end{tabular}
\end{center}
\step{2}{Tìm $b$}
$4 + 3b = 19 \Rightarrow 3b = 15 \Rightarrow b = 5$. Số thứ hai là $5$.
\step{3}{Viết tiếp đến số thứ tám}
$2,\ 5,\ 7,\ 12,\ 19,\ 31,\ 50,\ 81$ \quad ($12 + 19 = 31$; \ $19 + 31 = 50$; \ $31 + 50 = 81$). Số thứ tám là $81$.
\step{4}{Trung bình cộng của năm số đầu}
$(2 + 5 + 7 + 12 + 19) : 5 = 45 : 5 = 9$.
\step{5}{Quy luật chẵn -- lẻ của dãy Fibonacci}
\begin{recallbox}lẻ $+$ lẻ $=$ chẵn; \ lẻ $+$ chẵn $=$ lẻ; \ chẵn $+$ lẻ $=$ lẻ.\end{recallbox}
Viết L (lẻ), C (chẵn) cho các số $1, 1, 2, 3, 5, 8, 13, 21, 34, \ldots$:
\begin{center}
\begin{tabular}{|c|c|c|c|c|c|c|c|c|c|}
\hline
\rowcolor{tmBlue!10} Thứ tự & 1 & 2 & 3 & 4 & 5 & 6 & 7 & 8 & 9\\ \hline
Số & 1 & 1 & 2 & 3 & 5 & 8 & 13 & 21 & 34\\ \hline
Chẵn/lẻ & L & L & \textbf{C} & L & L & \textbf{C} & L & L & \textbf{C}\\ \hline
\end{tabular}
\end{center}
Mẫu ``L, L, C'' lặp lại mãi mãi.
\begin{proofbox}Giả sử tới một chỗ nào đó ta có hai số liên tiếp là L, L. Khi đó số tiếp theo là L $+$ L $=$ C; số sau nữa là
L $+$ C $=$ L; số sau nữa là C $+$ L $=$ L. Ta lại được hai số liên tiếp L, L, nên chu kì ``L, L, C'' bắt đầu lại. Dãy mở đầu
bằng $1, 1$ (L, L), nên mẫu này lặp lại từ đầu đến cuối: các số ở vị trí chia hết cho $3$ là số chẵn, các số còn lại là số lẻ.\end{proofbox}
\step{6}{Xét vị trí $2026$}
$2026 = 3\times 675 + 1$, tức $2026$ chia cho $3$ dư $1$. Vị trí này ứng với chữ ``L'' đầu tiên trong một chu kì. Vậy số thứ
$2026$ là \textbf{số lẻ}.
\begin{warnbox}Không thể kết luận ``$2026$ là số chẵn nên số thứ $2026$ cũng chẵn''. Tính chẵn -- lẻ của \emph{vị trí} khác
với tính chẵn -- lẻ của \emph{số} ở vị trí đó; ví dụ số thứ $2$ là $1$ (lẻ).\end{warnbox}
\verify{Câu a): $2 + 5 = 7$, $5 + 7 = 12$, $7 + 12 = 19$ -- đúng số thứ năm là $19$. Câu c): kiểm tra mẫu trên các vị trí
nhỏ: vị trí $4 = 3\times 1 + 1$ là số $3$ (lẻ); vị trí $7 = 3\times 2 + 1$ là số $13$ (lẻ).}
\ans{a) số thứ hai là $5$, số thứ tám là $81$;\quad b) $9$;\quad c) số lẻ}
}
